% !TeX root = ../papers/nc-p3.tex
\section{二つの二直線列と一つの三次列の五重接触}
\label{sec:cubic-corridor-jets}
\begin{theorem}\label{thm:cubic-corridor-jets}
純関係の有無によらず、順中心を位置 $0,1,2$ に一つずつ持ち、そこから始まる列が順に
$(1,1),(1,1),(1,3,1)$ である標準正則窓は存在しない。
\end{theorem}
\begin{proof}
位置0のねじれ三次曲線上の順・逆欠損は次数公式により各5で、それぞれ相異なる唯一の点
$F_0,B_0$ に集中する。位置 $j=1,2,3$ の初回・終回直線を $L_j,M_j$ と書く。
これらはそれぞれ一点で交わる。実際 $L_1$ には $F_1$、$M_2$ には $B_2$ があり、
$F_1$ の後続直線 $L_2$ は $L_1$ からの対応点でも $M_2$ と会う。
他の二組は三次曲線上の $F_0,B_0$ の後点・前点で会う。
初回と終回が一致すれば移動点一つの列となるので異なる。
定理\ref{thm:final-pure-branch-frames}の三中心の座標を使い、
三次を $C(1,x)=(1,x,x^2,x^3)$、$F_0,B_0$ を $0,\infty$ とする。
$C$ の後点は $M_1(1,x)=(0,1,x,0)$、前点は $L_3(1,x)=(1,x,0,0)$ である。
以下 $v_j$ を位置1の四つの座標関数、$y_j$ を位置3の四つの座標関数とする。

$R_0$ の $L_1$ への制限は完全な六次元である。よって三次上でその関係はちょうど
\[
 p_1h+p_2k+v_3D(h,k),\qquad
 (p_1,p_2)=(xv_0,xv_1-v_2),\qquad \deg h,\deg k\le2
\]
と表せる。実際 $v_0,v_1$ の係数は $F_0$ で消える任意の三次式であり、$v_2$ の係数は
後点から決まる。$D$ の値は次数3以下である。四つの評価行の三次小行列の共通因子は、
$h,k,D$ の評価三行の三次小行列の共通因子に原点でちょうど $x$ を掛けたものである。
$k$ の行は $v_2$ に単元で現れ、$h$ の行は $v_0$ に $x$ 倍で現れるからである。
従って $D(h,k)=[H(x)(h,k)^{\mathsf t}]_{<4}$ となる。
ここで $H$ は二成分の行、括弧は四次以上を捨てる操作である。実際 $h,k$ の定数基底列に
単位小行列があり、それらの列から読む $H$ で全列について $x^4$ を法として一致する。
原点で評価階数2なので $H(0)$ の第一成分は非零である。第四座標方向の倍率と
最初の三方向による平行移動で
\[
 H=H_0+x^2 H_2+x^3 H_3,\qquad H_0=(1,0)
\]
とできる。消した三係数は第一成分の一次係数、第二成分の定数と一次係数である。
$U=(H_2^{\mathsf t},H_3^{\mathsf t})^{\mathsf t}$（行は $H_2,H_3$）は非零である。
零なら三次上で一般評価階数が2に落ちる。
逆に $t=1/x$ で $R_3$ を読むと、座標 $0,2$ を交換した同じ表示になり、
列 $J(t)=J_0+t^2J_2+t^3J_3$、$J_0=(1,0)^{\mathsf t}$、$V=(J_2,J_3)\ne0$ を得る。

両辺の三文字重なりの元を $C(1,x)$ で評価する。$v_3,y_3$ を切除した部分は一意に
\[
 (p_1,p_2)(E_0+xE_1)q,\qquad q=(y_2,y_1-xy_0)^{\mathsf t}
\]
と書ける。実際 $h,k$ の $y_0$ の係数は、前点の消滅から $y_1$ の係数の $-x$ 倍なので、
$y_1$ の係数は一次以下である。$y_2$ の係数も一次以下である。
$B_0$ で $M_3$ を消すから $y_2$ の三次係数が零で、$p_1,p_2$ の一次係数は独立だからである。
この切除は重なりに単射である。一方の辺の三座標への制限は初回直線への単射を含み、
他辺も終回直線への単射を含むので、順に第四座標の行・列を復元できる。

$v_3$ の行を $A(x)q$ とすると、$A$ は二次以下でなければならない（前点の消滅と
$B_0$ での $y_2$ の三次係数の消滅）。そのための条件と係数は
\[
 H_2E_1+H_3E_0=0,\qquad (A_0,A_1,A_2)=(H_0E_0,H_0E_1,H_2E_0).
\]
逆側で $y_3$ の列を $(p_1,p_2)D'(x)$ とすると
\[
 E_0J_2+E_1J_3=0,\qquad (D'_0,D'_1,D'_2)=(E_1J_2,E_0J_0,E_1J_0).
\]
最後の $v_3y_3$ の係数は $[HD']_{<4}$ と $x^3[(t^2 A(1/t))J(t)]_{<4}|_{t=1/x}$ の
両方に等しい。その四係数を比較すると
\[
 H_2E_0J_3=H_2E_0J_2=H_2E_1J_2=H_3E_1J_2=0 .
\]

$U,V$ ともに可逆なら $X_j=UE_jV$ と置く。$X_0$ の第一行、$X_1$ の第一列が零で、
二つの線型条件が残りの行・列も零にする。$U$ が可逆で $V$ が階数1なら
$J_2=v\beta,J_3=v'\beta$ とする。$v\ne0$ なら最後の二条件から $E_1\beta=0$、
次いで $E_0\beta=0$。$v=0$ なら列条件から $E_1\beta=0$、混合条件と行条件から
$H_2E_0\beta=H_3E_0\beta=0$ で同じ結論となる。両行列が $\beta$ を消す四次元空間で
行条件は非自明なので重なり次元は高々3である。逆の場合も同じである。
いずれも位置2による独立な四縮約が重なりに入ることに反する。

残る階数1同士で $H_2=u\alpha,H_3=u'\alpha$ と書く。$uv\ne0$ なら混合条件から
$\alpha E_0\beta=\alpha E_1\beta=0$。$u=0$ なら行条件が $\alpha E_0=0$ を与え、
$v\ne0$ なら最後の混合条件、$v=0$ なら列条件から $\alpha E_1\beta=0$ を得る。
$v=0,u\ne0$ でも列条件と最初の混合条件を用いる。
従って $\alpha E_0,\alpha E_1$ はともに $\beta$ を消す一次元空間に属し、
$u\alpha E_1+u'\alpha E_0=0$ なので、その組の像次元は高々1である。
位置2の四縮約のうち三次元以上の部分空間で $E_0,E_1$ の像は固定直線
$k\lambda=\ker\alpha$ に入る。この部分では全ての係数が
\[
 x(\lambda_1v_0+\lambda_2v_1)+\lambda_1v_3-\lambda_2v_2
\]
で割れる。実際 $y_3$ も含め、対応する二次以下の列は $\lambda$ の倍であり、
$H_2\lambda=H_3\lambda=0$ なので第四行の復元に切捨てがない。
この一次式の $x$ の両係数を消す非零の位置1の点を取る。三次曲線は位置0全体を張るから、
この点で縮約した窓は位置2の三次元以上の点空間を消し、隣接支持が高々1になる。
標準正則窓では隣接評価階数は常に2以上であり、矛盾する。
\end{proof}
